\documentclass[10pt,a4paper]{amsart}
\usepackage[margin=19mm]{geometry}
\usepackage{amsmath,amssymb,mathtools}
\usepackage[scr=rsfs,cal=euler]{mathalfa}
\usepackage{xfrac}
\usepackage[unicode,hidelinks,bookmarks=false]{hyperref}
\newcommand{\ie}{i.e.\ }
\newcommand{\cf}{cf.\ }
\newcommand{\resp}{respectively\ }
\newcommand{\Subsection}{\subsection}
\providecommand{\nobreakdash}{\nobreak\hbox{-}}
\setlength{\emergencystretch}{2em}
\multlinegap0pt
\usepackage{amssymb}
\usepackage{mathtools}
\usepackage{xcolor}
\usepackage{tikz}
\usepackage[shortlabels]{enumitem}
\newtheorem{lem}{Lemma}
\title{\bf Mean Curvature Flow for Isoparametric Submanifolds in Hyperbolic Spaces}
\author{Xiaobo Liu, Wanxu Yang}
\date{Source: arXiv:2504.15602v3 (CC BY 4.0)}
\begin{document}
\maketitle

\noindent\emph{Source and licence.} \bf Mean Curvature Flow for Isoparametric Submanifolds in Hyperbolic Spaces. Version arXiv:2504.15602v3 (CC BY 4.0), distributed by arXiv under the Creative Commons Attribution 4.0 International licence, http://creativecommons.org/licenses/by/4.0/, as recorded in the arXiv licence metadata for this exact version and shown on the abstract page https://arxiv.org/abs/2504.15602v3. Full licence notice: \texttt{licenses/arxiv-2504.15602-CC-BY-4.0.txt}.\par\smallskip

\noindent\textit{Editorial scope.} Counted unit: the article's lem F'formula (source0.tex lines 1009--1025). The article's complete original proof of it is delivered with it (source0.tex lines 1027--1063, one \texttt{proof} environment of 37 lines). No mathematics is written, repaired or re-derived: the statement and the proof are the article's own lines, in the article's own notation, and no retained line is substituted or re-flowed.  Retained uncounted as the source-local context the proof needs: lines 530--535; lines 566--573; lines 837--839; lines 841--843; lines 869--879; lines 875--877.  Nothing was omitted from any retained span, so the package carries no omission record.  The retained lines cite 1 works, whose entries are transcribed verbatim from the article's own bibliography source in the e-print and delivered with short editorial labels recorded in the item's reference mapping.  No retained line contains a figure environment, an image-inclusion command or drawing code, so no image is delivered and the package declares no asset dependency.  Editorial formatting only, in addition: an AMS \texttt{amsart} preamble that restates the article's own declarations and macros used by the retained lines, namely the environments \texttt{lem} on separate counters, together with the standard packages those lines need.  The retained environments print in this document as Lemma 1 in order of appearance, whereas the article prints the same items with the numbering of its own full document.  The complete native source of the article as distributed in the arXiv e-print for this exact version is bundled unchanged as \texttt{sources/176880/source0.tex}.
\begin{lem} \label{4.1}
		Assume $M$ is an $n$-dimensional submanifold of $\mathbb{H}^{m}(-r) \subset \mathbb R^{m,1}$. Let $H(x)$ and  $H^L(x)$  be the mean curvature vectors of $M$ in $\mathbb{H}^{m}(-r)$ and  $\mathbb R^{m,1}$ respectively at $x \in M$. Then
	\begin{align*}
		H^L(x) = H(x) + \frac{n}{r}x.
	\end{align*}
\end{lem}

\begin{lem} \label{1.2}
		Let  $f(\cdot,t)$ and $F(\cdot,t)$ be respectively the hyperbolic MCF and Lorentzian MCF of an $n$-dimensional submanifold
$M$ in $\mathbb{H}^{m}(-r) \subset \mathbb R^{m,1}$.  Then for any $x \in M$,
		\begin{align}
			F(x,t)=\sqrt{1+\frac{2nt}{r}}f(x,\frac{r}{2n}\ln(1+\frac{2nt}{r})) \label{7.1.4}
		\end{align}
if both sides of this equation are well defined.
\end{lem}

\begin{equation} \label{eqn:condxi}
 a := \left<u, \, \xi \right> \geq 0, \hspace{40pt}  \langle\xi, \,\, \xi\rangle = \pm 1 \,\,\, {\rm or} \,\,\, 0.
\end{equation}

\begin{align} \label{eqn:V+u}
	 		V+u=\{v\in\mathbb R^{m,1}\mid \left<v,\xi \right>=a\}.
\end{align}

\begin{lem}\label{4.4}
	Assume $M$ is an $n$-dimensional submanifold of $L(V,u)  \subset \mathbb{H}^{m}(-1)$. Let $H_1(x)$ (respectively $H(x)$) be the mean curvature vector of $M$ in  $L(V,u) $ (respectively in $\mathbb{H}^{m}(-1)$) at $x\in M$. Then we have
	\begin{align*}
		H(x)= H_1(x) - n\alpha(\alpha x+\beta \xi),
	\end{align*}
	where
    \begin{equation} \label{4.12}
		\beta :=\frac{1}{\sqrt{\left<\xi, \, \xi \right>+a^2}}, \hspace{10pt}  \,\,\,\,\,\, \alpha :=\beta a,
    \end{equation}
    and $\xi$, $a$ are defined before equation \eqref{eqn:V+u}.
\end{lem}

    \begin{equation} 
		\beta :=\frac{1}{\sqrt{\left<\xi, \, \xi \right>+a^2}}, \hspace{10pt}  \,\,\,\,\,\, \alpha :=\beta a,
    \end{equation}

\begin{lem} \label{F'formula}
Assume $M$  is an $n$-dimensional  submanifold in a totally umbilical submanifold $L(V, u) \subset \mathbb H^{m}(-1)$.
Let $f_1(\cdot,t)$ and $F(\cdot,t)$ be the solutions to  mean curvature flows for $M$ in $L(V,u)$ and  $\mathbb R^{m,1}$ respectively. Let $\xi$, $\alpha$ and $\beta$  be defined by equations \eqref{eqn:condxi} and \eqref{4.12} .
If $\alpha\neq 1$, then for $x\in M$,
	\begin{align}
		F(x,t)=\sqrt{2nt(1-\alpha^2)+1}\, \, (f_1(x,s_{\alpha}(t))-\eta)+\eta,  \label{9.1.5}
	\end{align}
	where $\eta := \frac{\alpha\beta}{1-\alpha^2}\xi$ and  $s_{\alpha}(t)$ is defined by
	\begin{align} \label{eqn:st}
		 	s_{\alpha}(t) :=
		 	\frac{1}{2n(1-\alpha^2)}\ln (2nt(1-\alpha^2)+1).
     \end{align}
If $\alpha=1$, then
\begin{align}
		  	 F(x,t)&=f_1(x,t)-  nt  \beta \xi .  \label{9.1.6}
		  \end{align}
\end{lem}

\begin{proof}
If $\alpha \neq 1$, it is straightforward to check that $\eta \in V + u$ and $L(V, u)$ consists of points $x \in V+u$ satisfying equation
$ \langle x-\eta, x-\eta\rangle = (\alpha^2-1)^{-1}. $
Let $\tau$ be the translation on $\mathbb R^{m,1}$ defined by $\tau(x)=x-\eta$ for $x\in \mathbb R^{m,1}$.
Then $\tau(L(V, u))$ is a sphere (or hyperbolic space)  \footnote{$\tau(L(V, u))$ is a circle or hyperbola if its dimension is one.}
%with sectional curvature $\alpha^2-1$
centered at origin in the Eulidean (or Lorentz) space
$\tau(V+u)$ if $\alpha > 1$ (or $0 \leq \alpha < 1$).
Since $\tau$ is an isometry, $\tau(f_{1}(\cdot,t))=f_{1}(\cdot,t)-\eta$ is the solution to MCF for $\tau(M)$ in $\tau(L(V,u))$. By Lemma \ref{1.2}  and its spherical analogue (cf. Equation (1.1) in \cite{X.C}), the function
\[ G(x, t) := \sqrt{2nt(1-\alpha^2)+1}\, \, (f_1(x,s_{\alpha}(t))-\eta)\]
 is the solution to MCF for $\tau(M)$ in $\tau(V+u)$. Since $\tau(V+u)$ is a totally geodesic submanifold in $\mathbb R^{m,1}$. $G(x, t)$ also solves the MCF of $\tau(M)$ in $\mathbb R^{m,1}$.
 Hence the solution of MCF for $M$ in $\mathbb R^{m,1}$ is $\tau^{-1}(G(x,t))$, which is exactly the
 function given by equation \eqref{9.1.5}.

If $\alpha=1$, then $\langle \xi, \xi \rangle =0$ and the induced metric on $\tau(V+u)$ is degenerate. In this case, MCF for submanifolds in $\tau(V+u)$ is not well defined. So the above arguments do not work for this case.
Instead, we will check directly that the function $F(x, t)$ in equation \eqref{9.1.6}, which is the limit of the function in
equation \eqref{9.1.5} as $\alpha \rightarrow 1$, satisfy the MCF equation.
Let $H_1(\cdot,t)$ be the mean curvature vector field of $f_1(\cdot,t)$ in $L(V,u)$. Then
\begin{align}
\frac{\partial}{\partial t}F(x,t)
&=H_{1}(x,t)-n\beta\xi.\label{2.9}
\end{align}
On the other hand, since $f_1(\cdot ,t)\subset L(V,u) \subset \mathbb{H}^{m}(-1)$,
by Lemma \ref{4.4}, the mean curvature vector of $f_1(\cdot ,t)$ in $\mathbb{H}^{m}(-1)$ at point $f_1(x,t)$ is
		 \begin{align*}
		 	H_1(x,t) -n(f_1(x,t)+\beta \xi).
		 \end{align*}
By Lemma \ref{4.1},
the mean curvature vector of $f_1(\cdot ,t)$ in $\mathbb R^{m,1}$ at point $f_1(x,t)$ is
$ H_1(x,t)- n \beta \xi$,
which is exactly the right hand side of equation \eqref{2.9}.
Since for each fixed $t$, $F(\cdot, t)$ is just a translation of $f_{1}(\cdot, t)$,
the right hand side of equation~\eqref{2.9} is also the mean curvature vector of
$F(\cdot,t)$  in $\mathbb R^{m,1}$ at the point $F(x,t)$.
Hence $F(\cdot, t)$ is the solution to MCF of $M$ in $\mathbb R^{m,1}$.
This completes the proof of the lemma.
\end{proof}

\begin{thebibliography}{99}
\bibitem[X.C]{X.C} Liu, X. and Terng, C.L., {\it Ancient solutions to mean curvature flow for isoparametric submanifolds}, Math. Ann. Vol 378(2020), pp 289-315.
\end{thebibliography}
\end{document}
